% Peplinski (2026), SAFE profile validation. Compile with pdflatex twice; the bibliography is inline (no BibTeX).
% Code, frozen analysis plan and aggregate results: https://doi.org/10.5281/zenodo.24396236
\documentclass[11pt]{article}
\usepackage[utf8]{inputenc}
\usepackage[T1]{fontenc}
\usepackage[a4paper,margin=2.5cm]{geometry}
\usepackage{amsmath,amssymb}
\usepackage{booktabs}
\usepackage{threeparttable}
\usepackage{array}
\usepackage{microtype}
\usepackage[hidelinks]{hyperref}
\usepackage{xcolor}
\usepackage[round,authoryear]{natbib}

\newcommand{\nats}{\,\text{nats}}

\title{Does a Firm Profile Help an AI Agent Forecast Real Survey Answers?\\
\large A Prospectively Specified Validation on the EC/ECB Survey on the Access to Finance of Enterprises for Polish Firms}
\author{Szymon P. Pepliński\\ \small Generatywnie\\ \small ORCID 0009-0008-2928-0655}
\date{20 November 2026}

\begin{document}
\maketitle

\begin{abstract}
Language models are increasingly used as simulated respondents, yet it is rarely tested whether the detailed
profile given to such an agent carries checkable information about how the real unit answers, beyond what a coarse
segment already implies. We test this with firm-level records of the EC/ECB Survey on the access to finance of
enterprises (SAFE) for Poland. For each firm, three language models receive either a full profile card built from
other answers of the same interview (size, sector, turnover, age, ownership, exports, which sources of financing are
relevant and were used, and financing applications with their outcomes or the stated reason for not applying) or a
short card with size and sector only, and give a probability distribution over the firm's answer to how important access
to finance has been as a problem (1--10). The primary measure is the band log score of the observed answer; the
primary comparison per model is full versus short card, tested with a directional studentised bootstrap and Holm's
procedure across the three models. After three development pilots, the analysis plan was frozen before model
evaluation on the 2026Q3 wave (692 respondents). Aggregate test-answer counts had been seen during development;
this exposure is disclosed separately from the subsequent confirmatory evaluation. For one model the full profile improves
the forecast over the short card by 0.059\nats{} (95\% CI 0.022 to 0.095; Holm-adjusted rejection); for a second no
improvement is detected at the precision reached (individual 95\% CI $-0.021$ to $0.062$\nats); for the third the full card is worse.
However, on the primary band log score the forecasts of all three models, with either card, score below the plain
marginal distribution of answers, and a penalised logistic regression on the same card fields outperforms them. The profile thus helps a weak segment
forecast within a model, without establishing an advantage over simple statistical baselines.
\end{abstract}

\section{Introduction}

Large language models are used to simulate survey respondents and economic agents
\citep{argyle2023,horton2023,aher2023}, and the question of whose views such simulations reproduce has received
attention \citep{santurkar2023}. In applied settings, an agent is often given a rich profile of the unit it stands
for. Whether that profile carries information that can be checked against the unit's real answers, beyond a coarse
segmentation such as firm size and sector, is an empirical question that a simulation alone cannot settle.

This study belongs to a project that builds a population of Polish firms from public data and uses language-model
agents to forecast how firms in that population would respond. The premise of such a framework is that a well-built
firm profile carries information about the firm's real answers that a coarse segment does not. We test that premise
where it can be checked against real respondents: in SAFE, where the profile and the answer come from the same firm.
The profile tested here is therefore built from other answers of the same SAFE interview, not from public
registers: structural characteristics, the use of financing, applications and their outcomes, and stated reasons for
not applying. Such a profile may carry more information about the answer than a register-based one, so the result
does not show that a register-based profile would work as well. The SAFE terms rule out any record-level linkage with
register data, so a register-based profile cannot be scored against SAFE answers firm by firm. Comparisons of the simulated population with survey distributions are a separate,
distribution-level step not reported here.

We ask: \emph{does a firm's profile improve a language model's forecast of that firm's real survey answer over a
forecast from its size and sector alone?} We answer it on public-sector survey microdata with a design that separates
development from confirmation: the elicitation, the scoring, the tests and their calibration were fixed in a plan
that was frozen before any model saw the held-out wave. A positive answer for a model would support testing richer
agent profiles in later product experiments; it would not replace them. No product or offer was shown to the
respondents, so the study does not validate product choices, purchase intentions or product decisions.

\section{Data}

\paragraph{Survey.} We use the Polish records of SAFE, a survey of enterprises run by the European Central Bank and
the European Commission, from the waves of 2024Q3, 2025Q3 and 2026Q3, accessed under the ECB's SAFE microdata user
agreement for scientific use. The target question is item 3 of question Q0b: ``How important have the following
problems been for your enterprise during the previous quarter and current quarter? \ldots{} Access to finance'', on a
scale from 1 (not at all important) to 10 (extremely important). Answers ``don't know'' are excluded from scoring and
counted. The target question was chosen after a screening of several candidate questions on the development data only:
for each, a penalised multinomial logit on the candidate card fields was compared with the marginal distribution in
cross-validation grouped by firm, and Q0b\_3 showed one of the larger gains (0.060\nats{} in band log-loss, against
0.021 for a question on late payments). The choice therefore rests on a signal found in development data; it was made
before any model saw the test wave, and the frozen plan records it.

\paragraph{Split.} The \emph{test} set is the 2026Q3 wave: 722 records, of which 692 answer Q0b\_3 with 1--10. The
\emph{development} set contains 981 eligible firm--wave observations from 952 distinct firms in the 2024Q3 and
2025Q3 waves, excluding every firm present in the test wave. Development pilots sampled one record per firm;
the statistical baselines for the confirmatory analysis were trained on all 981 observations.

\paragraph{Information available before freezing.} The test wave was reserved for confirmatory model evaluation,
but its aggregate answers were not entirely unseen. The frozen plan records that reviewers (AI review agents; see Use of AI tools) had viewed the Q0b\_3
marginal counts, including the 30 excluded answers, and aggregates for a separate item, qa21. The plan states that
no recommendation was based on the exposed Q0b\_3 marginal. Test-wave size-class shares and the eligible sample size
were used for planning precision and calibration; firm identifiers were used to exclude overlap with development
data. These aggregate exposures limit any claim of a completely untouched holdout. The elicitation and analysis
procedure were frozen before the confirmatory model evaluation.

\paragraph{Profile cards.} The \emph{full card} states the firm's employment class, sector (four groups), annual
turnover class, age class, the largest owner, the export share, for eleven sources of financing whether each is
relevant to the firm (used in the past or considered for the future) and, if so, whether it was used in the
period, and for up to four kinds of financing whether the firm applied and with what outcome or, if it did not
apply, the stated reason (fear of rejection, sufficient internal funds, or other reasons). The card is thus a mix of
structural characteristics, financing events and self-reported intentions and motives. Some of these are close in
content to the target, a self-assessment of access to finance as a problem: in particular, not applying for fear of
rejection, and not needing external funds. All of them were recorded in the same interview as the answer being
forecast, so the task is to reconstruct a concurrent answer, not to forecast later behaviour. The card excludes the
questions that the card key classes as perceptions or expectations (all items of Q0b, among them the target, and
Q2, Q11, Q16, Q17 and Q19), the follow-up questions on late payments, and any identifier, weight or period. The \emph{short card} states employment class
and sector only. A \emph{shuffled card} keeps the firm's size and sector and takes the remaining blocks from a donor
firm of the same size$\times$sector cell; it is a diagnostic of whether matching the other features to the firm
matters. Cards, the key that explains every code and the question are given to the models in Polish; the Polish
wording is our translation of the English questionnaire.

\section{Elicitation and models}

\paragraph{Elicitation.} Each card is a separate conversation. The system prompt states the question (with the
context of the Q0b battery), the key to the card and the task: imagine 100 firms drawn at random from firms in Poland
that share every feature on the card, and state how many of them would choose each rating from 1 to 10; ratings that
cannot be ruled out must not receive zero. The integers are normalised to a distribution. This form (``arm B'') was
chosen on development data from three alternatives by a rule fixed beforehand (Section~\ref{sec:dev}).

\paragraph{Models.} Three models were run with fixed settings, the same in the last development pilot and in the test:
Qwen~3.8 27B served locally by Ollama (temperature 0, fixed seed, thinking off), Claude Haiku~5.5 through the
Anthropic Messages API (effort ``medium'', thinking disabled) and GPT-6 Luna through the OpenAI Responses API
(reasoning effort ``none''). The API models are identified by the model name the API reported, which may be an alias;
the analysis compares, for every session, the SDK and API versions and the model's reported creation date with the
frozen values.

\paragraph{Repeated draws.} Model noise in the comparison was estimated on development data from repeated draws of
both cards for the same firms. By a rule fixed in advance (the 80\% upper bound of the noise share of the variance of
the per-firm difference below 0.2: one draw; below 0.5: two; else three), the confirmatory analysis averages the band
log scores of $K=1$ draw for Qwen and Haiku and $K=2$ draws for Luna. This rule applies to the full and short cards.
The shuffled-card diagnostic uses one shuffled draw per firm for every model; for Luna it therefore compares the
average of two full-card log scores with one shuffled-card log score.

\section{Scoring and inference}
\label{sec:inference}

\paragraph{Score.} Each forecast $p$ over the ten ratings is smoothed towards the uniform distribution,
$\tilde p_j=0.99\,p_j+0.001$, and collapsed to three bands (1--3, 4--7, 8--10). The per-firm
\emph{band log score} is the logarithm of the resulting probability of the observed band; higher is better.
The ordinary logarithmic score is strictly proper for the distribution being evaluated \citep{gneiting2007}.
Here the evaluated distribution is the smoothed three-band forecast. We do not claim strict propriety for the
entire transformation as a scoring rule on the raw ten-rating forecast: collapsing bands discards within-band
information, and smoothing changes the probabilities before scoring. For firm $i$ and a model,
\begin{equation}
d_i = S_i(\text{full}) - S_i(\text{short}),
\end{equation}
averaged over the $K$ draws of each card. The ten-value log score and the ranked probability score \citep{epstein1969}
are reported as explanatory measures.

\paragraph{Primary family.} For each model, $H_1$: the mean of $d$ is not positive. Each is tested with a two-sided
studentised bootstrap (bootstrap-$t$, 10{,}000 resamples of firms; \citealp{hall1992,efron1993}) and a directional
rule ($p$ if $\bar d>0$, otherwise 1). The family of the three $H_1$ is controlled by Holm's procedure at $\alpha=0.05$
\citep{holm1979}. A rejection supports the claim for that model only. The bootstrap resamples firms together with the
scores their forecasts obtained; it does not call the models again. The intervals therefore describe the uncertainty
of a mean over respondents drawn like these, given the answers the models gave, and include model noise only through
the realised per-firm scores; they do not measure the variation of repeated model answers on the same cards alone.
Nor do they support inference to all Polish firms: the respondents are an unweighted set of complete cases from one
survey wave.

\paragraph{Valid answers and retries.} An answer is valid when it is one JSON object with exactly two keys, the counts
for the ratings 1 to 10 (each an integer from 0 to 100, not all zero) and a text reason; duplicate keys, other keys,
non-integers and text around the object make it invalid. The counts are divided by their sum, which need not be 100.
The instruction not to give zero to ratings that cannot be ruled out is not enforced by the parser; zeros are handled
by the smoothing of the score. A refusal, a cut-off or an invalid answer is asked again in a fresh conversation, up to
two attempts in all; the first valid answer is used, and an item still without one is a technical error, never a
survey answer. A retry repeats the same item and is not one of the $K$ draws: the $K$ draws are separate items of the
build, each with its own up to two attempts. For the local model at temperature 0 with a fixed seed, a retry usually
repeats the invalid answer.

\paragraph{Estimand and missing answers.} The estimand is the mean $d$ among the respondents for whom all $K$ draws of
both cards gave a valid answer (complete cases). Coverage is reported with three sensitivities over all eligible
respondents: invalid answers scored as the marginal distribution, scored worst in both arms, and scored against the
claim (worst for the full card, marginal for the short card).

\paragraph{Mandatory references and exploratory analyses.} For every model the differences full minus marginal and
short minus marginal are reported, together with the diagnostic full minus shuffled. An exploratory comparison
($H_3$) of the full card with a penalised multinomial logit on size and sector, and statistical baselines, are reported
without any confirmatory claim. All statistical baselines are estimated on the 981 development observations and
predict the ten ratings; their forecasts are smoothed and scored like the models'. The \emph{marginal distribution}
is the distribution of the ten ratings in the development data with 0.5 added to each of the ten counts. The
\emph{logit} is an L2-penalised multinomial logit on the 25 card columns (six structural fields, eleven sources of
financing, four applications and four outcomes), with the penalty chosen by four-fold cross-validation grouped by firm
on the same smoothed band log score; the logit on size and sector uses the same procedure on those two fields, and an
explanatory ablation leaves out the applications and their outcomes. Each comparison uses the firms with valid answers
for the cards it involves, so denominators differ between comparisons (Table~\ref{tab:ref}). The intervals for
differences from the marginal and for $H_3$ are symmetric studentised bootstrap intervals; those for full minus
shuffled are percentile bootstrap intervals.

\paragraph{Calibration of the test.} Before freezing, the level of the directional bootstrap-$t$ was checked on each
model's own centred development differences, at the expected number of complete cases and at 692 firms, for $H_1$ and
$H_3$, at every local level a graphical procedure over $H_1$ and $H_3$ could reach (0.0167, 0.025, 0.05): 4{,}000 outer
replications with 10{,}000 inner resamples, and simultaneous one-sided 95\% Clopper--Pearson bounds \citep{clopper1934}
with a Bonferroni correction over all 36 bounds. The outer replications resample the empirical distribution of the
centred development differences, so the check concerns the level under that distribution, not under every
distribution the test sample could have. All 18 bounds for $H_1$ lay within their levels; four bounds for $H_3$
(two models) did not, with observed rates below their levels. Because a graphical procedure \citep{bretz2009,bretz2011}
would pass the level of a rejected $H_3$ on to the other $H_1$, the plan's open choice between that procedure and Holm
over the three $H_1$ was resolved for Holm before the confirmatory analysis, and $H_3$ became exploratory. This decision
was taken after development results and is reported as such.

\section{Development stage}
\label{sec:dev}

Three development pilots and one comparison of elicitations preceded the freezing; they informed the design and are
not evidence for the claim. Table~\ref{tab:dev} summarises the mean $d$.

\begin{table}[ht]
\centering
\begin{threeparttable}
\caption{Development pilots: mean $d$ (full minus short card, band log score, nats).}
\label{tab:dev}
\small
\begin{tabular}{llp{6.0cm}rrr}
\toprule
Pilot & Sample & Elicitation & Qwen & Haiku & Luna\\
\midrule
1 & 120 & probabilities, no zero for ratings that cannot be ruled out & $+0.100$ & $+0.167$ & $+0.067$\\
2 & 120 & counts per 100, zeros allowed & $-0.302$ & $-0.086$ & $+0.003$\\
3 & 101 & counts per 100, arm B (chosen) & $-0.232$ & $-0.022$ & $+0.088$\tnote{a}\\
\bottomrule
\end{tabular}
\begin{tablenotes}\footnotesize
\item Sample is the number of firms selected, not the complete-case denominator. In Qwen/Haiku/Luna order,
the numbers contributing to the reported means are 84/117/120 in pilot 1, 112/120/120 in pilot 2, and 100/101/101
in pilot 3. Differences across pilots reflect changes in firms, elicitation and complete-case coverage.
\item[a] With $K=2$ draws; $+0.076$ with the first draw. All other entries use the first draw.
Pilot 3 selected 30 firms in each of the three smaller size classes and all 11 remaining firms with 250 or more employees.
\end{tablenotes}
\end{threeparttable}
\end{table}

\paragraph{Diagnosis after pilot 2.} In pilot 2 Qwen's full-card answers often put zero on the whole band 8--10, and
nine of 112 firms whose observed answer had near-zero probability accounted for 81\% of the negative sum of $d$; its
technical errors were valid objects with every count zero. The causal role of the instruction was a hypothesis, tested
next.

\paragraph{Comparison of elicitations.} On the same 120 firms and the same cards, three instructions were compared:
A, counts per 100 with the zero rule of pilot 2; B, the same with only the zero rule replaced by that of pilot 1; C, the
probabilities instruction of pilot 1. The rule, fixed before the calls, ranked arms by the mean band log score relative
to the marginal with equal weights of models and cards, excluded arms with more than 5\% technical errors for any model
and card, excluded an eligible arm only when another eligible arm had a demonstrated advantage (paired bootstrap over
firms, family-wise 95\% over the three pairs), and among the remaining arms chose the smallest change (A, then B, then
C) --- a conservative choice under an unresolved difference, not a claim of equivalence \citep{lakens2017}. The mean
scores were $-0.105$ (A), $-0.086$ (B) and $-0.075$ (C) with no demonstrated difference; only B met the error limit,
and B was chosen.

\paragraph{Pilot 3 and freezing.} Arm B was then run on 101 new development firms. Technical requirements were met
(one Qwen full-card first draw was invalid; no short-card first draws were invalid), and the calibration check of
Section~\ref{sec:inference} was run on these data. The frozen settings bind the confirmatory analysis to the very test
set (SHA-256 of the source, items and key), to the run settings and $K$ of each model, to the SHA-256 of the scripts,
the card key and the plan, and to the runner modules and serving stack recorded in every session.

\section{Results}

\paragraph{Primary family.} Table~\ref{tab:main} gives the confirmatory result on the 2026Q3 wave.

\begin{table}[ht]
\centering
\begin{threeparttable}
\caption{Confirmatory comparison on the 2026Q3 test wave (692 respondents): full versus short card.}
\label{tab:main}
\small
\setlength{\tabcolsep}{3.5pt}
\begin{tabular*}{\textwidth}{@{\extracolsep{\fill}}lcccccc@{}}
\toprule
Model & $K$ & Firms (coverage) & Mean $d$ (nats) & 95\% CI & Directional $p$ & Holm\\
\midrule
GPT-6 Luna & 2 & 692 (1.00) & $+0.059$ & $[0.022,\ 0.095]$ & 0.0019 & rejected\\
Claude Haiku 5.5 & 1 & 689 (0.996) & $+0.021$ & $[-0.021,\ 0.062]$ & 0.32 & not rejected\\
Qwen 3.8 27B & 1 & 685 (0.990) & $-0.192$ & $[-0.267,\ -0.116]$ & 1 & not rejected\\
\bottomrule
\end{tabular*}
\begin{tablenotes}\footnotesize
\item Individual symmetric studentised bootstrap 95\% intervals, not simultaneous intervals for the family.
Holm adjusts the hypothesis-test decisions. Luna had no technical errors, so all sensitivities equal the main
result; the sign-flip and $t$ tests give $p=0.0016$ and $p=0.0019$. For Haiku, scoring invalid answers worst gives
$+0.0003$. Qwen's errors were differential (7 full cards, 0 short cards); the sensitivity against the claim is
$-0.221$.
\end{tablenotes}
\end{threeparttable}
\end{table}

For Luna, the full profile improves the band log score over the short card by 0.059\nats{} on average among complete
cases, and the rejection survives Holm's correction across the three models. For Haiku, no improvement is detected at
the precision reached: the estimate is $+0.021$\nats{}, with an individual 95\% interval from $-0.021$ to
$+0.062$\nats{}. Its upper endpoint describes the uncertainty under this procedure; it is not a certain bound
or a demonstration of practical equivalence. For Qwen, the full card lowers the score.

\paragraph{References.} Table~\ref{tab:ref} reports the mandatory references.

\begin{table}[ht]
\centering
\begin{threeparttable}
\caption{Band log score relative to references (nats; individual descriptive 95\% CI), test wave.}
\label{tab:ref}
\footnotesize
\setlength{\tabcolsep}{3pt}
\begin{tabular}{lcrcrcr}
\toprule
Model & Full minus marginal & $n$ & Short minus marginal & $n$ & Full minus shuffled & $n$\\
\midrule
GPT-6 Luna & $-0.062\ [-0.117,\ -0.008]$ & 692 & $-0.121\ [-0.164,\ -0.077]$ & 692 & $+0.186\ [0.137,\ 0.238]$ & 692\\
Claude Haiku 5.5 & $-0.050\ [-0.088,\ -0.013]$ & 689 & $-0.070\ [-0.096,\ -0.044]$ & 692 & $+0.052\ [0.001,\ 0.103]$ & 686\\
Qwen 3.8 27B & $-0.359\ [-0.455,\ -0.263]$ & 685 & $-0.179\ [-0.236,\ -0.121]$ & 692 & $+0.279\ [0.190,\ 0.368]$ & 681\\
\bottomrule
\end{tabular}
\begin{tablenotes}\footnotesize
\item $n$: firms in the comparison (valid answers for every card it involves). The first two columns use different
firms where $n$ differs, so their difference need not equal the $H_1$ estimate of Table~\ref{tab:main} (Qwen:
$-0.180$ against $-0.192$). Studentised intervals against the marginal, percentile intervals for full minus shuffled.
\end{tablenotes}
\end{threeparttable}
\end{table}

In the test wave (692 firms) the band log-loss of the marginal distribution is 1.061, of the logit on size and sector
1.060 and of the logit on all card fields 0.981, a gain of 0.081\nats{} over the marginal; without the applications
and their outcomes the logit reaches 1.004, between the two. The exploratory $H_3$ differences (full card minus the
logit on size and sector, on the complete cases of the full card: 692, 689 and 685 firms) are $-0.064$
$[-0.118, -0.010]$ for Luna, $-0.052$ $[-0.090, -0.014]$ for Haiku and $-0.361$ $[-0.456, -0.265]$ for Qwen.

\paragraph{Explanatory measures.} Table~\ref{tab:expl} reports the two explanatory measures announced in the plan.
They are descriptive: no interval or test was planned for them, and none is reported. On the ten-value log score
every model forecast is worse than the marginal. On the ranked probability score Luna's full card (0.192) is
descriptively better than the marginal (0.198), on the same 692 firms; Haiku's and Qwen's full cards are not. The
statements below about losing to the marginal therefore refer to the primary band log score.

\begin{table}[ht]
\centering
\begin{threeparttable}
\caption{Explanatory measures on the test wave: ten-value log-loss and ranked probability score (lower is better).}
\label{tab:expl}
\small
\begin{tabular}{lrrrrr}
\toprule
 & \multicolumn{2}{c}{Ten-value log-loss} & \multicolumn{2}{c}{RPS} & \\
\cmidrule(lr){2-3}\cmidrule(lr){4-5}
Forecast & Full & Short & Full & Short & Firms\\
\midrule
GPT-6 Luna & 2.377 & 2.385 & 0.192 & 0.212 & 692\\
Claude Haiku 5.5 & 2.230 & 2.284 & 0.205 & 0.216 & 689\\
Qwen 3.8 27B & 2.587 & 2.453 & 0.218 & 0.211 & 685\\
\midrule
Marginal distribution & \multicolumn{2}{c}{2.059} & \multicolumn{2}{c}{0.198} & 692\\
Logit on size and sector & \multicolumn{2}{c}{2.055} & \multicolumn{2}{c}{0.198} & 692\\
Logit on all card fields & \multicolumn{2}{c}{1.963} & \multicolumn{2}{c}{0.175} & 692\\
\bottomrule
\end{tabular}
\begin{tablenotes}\footnotesize
\item Models: complete cases of the primary comparison; statistical forecasts: all 692 eligible firms. Ten-value
log-loss and RPS of the smoothed ten-rating forecast.
\end{tablenotes}
\end{threeparttable}
\end{table}

\section{Discussion}

The result for Luna is the configuration the plan anticipated as ambiguous: the full profile improves the model's
forecast over its own short-card forecast. The descriptive full-minus-shuffled comparison is positive, consistent
with a benefit from matching the other features to the firm. However, on the primary band log score both the
short-card and full-card forecasts score below the marginal distribution (on the explanatory RPS Luna's full card is
descriptively slightly better than the marginal; Table~\ref{tab:expl}). In the plan's words, the
profile helps a weak segment variant while the model loses to the plain distribution of answers. The statistical
baseline on the same fields, a regularised logit, beats every model's full-card forecast (by 0.143\nats{} for Luna,
0.131 for Haiku and 0.434 for Qwen). A positive $H_1$ is therefore
not evidence of an advantage over a statistical model, nor of value for product decisions; it establishes that, for
one model and one task, the profile carries information the model uses. The statistical baselines learn from 981
development answers, which the models do not see; the comparison is one of available procedures, not a bound on what
language models can do. The ablation of the logit suggests that the application block, including the stated reasons
for not applying, carries part of the signal; it does not identify what the models use.

The development stage shows how sensitive such results can be. The mean $d$ of one model differed by about
0.4\nats{} between two pilots that differed both in the firms and in the elicitation; within one sample, most of the
negative sum came from a few confident misses that a log score penalises heavily. On the same firms and cards,
replacing only the zero rule (A versus B) raised that model's full-card score relative to the marginal from $-0.376$ to
$-0.252$. The planned comparison did not establish a difference between arms; this does not establish equivalence
or, by itself, identify insufficient power as the explanation.

\paragraph{What this means for population-level profiles.} The result speaks to the premise, not to a particular
population: it shows that one model can use profile information of the kind a population of firms would carry, on a
task where the answer is real. It does not show that profiles assembled from registers are as informative as survey
fields of the same firm, nor that the forecasts are good in absolute terms; on the primary measure they are not
better than the marginal distribution. It also scores individual forecasts, which is a different target from
reproducing the distribution of answers within a segment: a forecast can match a segment's distribution without
distinguishing firms, or improve individual forecasts while the segment's distribution stays wrong. Variation between
a model's draws for the same card is model noise, not evidence that the model reproduces the heterogeneity of similar
firms. A population-level check against the distributions of SAFE answers by segment is the natural next step and needs
no record-level linkage. Using a statistical model for a firm's needs and a language model for its reaction to a
specific offer is a working hypothesis that this study motivates but does not test; testing it needs data in which
real firms reacted to described offers.

\paragraph{Limitations.} The claim is about complete cases of the 2026Q3 Polish SAFE respondents answering 1--10, not
the population of firms; weighted results are explanatory only (Luna, weighted: $+0.029$). One question and one country
are covered. The API models are identified by reported names that may be aliases. The calibration of $H_3$ was not
confirmed for two models, which is why $H_3$ is exploratory; the switch from the graphical procedure to Holm was made
after development results. The development pilots informed the design rather than serving as confirmatory evidence.
Aggregate test-answer counts had been seen before freezing, as disclosed in the Data section; the test wave should
not be described as a completely untouched holdout. The order of freezing and confirmatory runs is documented by the
frozen settings (with the approval date entered by the author), the checksums that bind the analysis to them, and the
commit history of the private project repository; it was not registered with an independent registry, and we do not
describe it as a preregistration. The plan asked for each model's stated training cutoff to be recorded before
freezing; this was not done. Record-level SAFE data are available under the ECB's user agreement and published
aggregates of earlier waves are public; whether a model had seen SAFE files or aggregates in training, including the
2026Q3 public-use file published in late October 2026, was not established. This is a limit on interpretation, not
evidence of memorisation, and it is separate from the researchers' disclosed exposure to aggregates.

\section*{Use of AI tools}

In addition to the three language models evaluated in this study, the author used large-language-model-based AI
assistants throughout the work to write most of the code according to the author's specifications, run analyses and
checks, and verify references. The author reviewed the code at each step and requested corrections where needed. The
code reviews, blind-agent checks of the study materials, and domain review rounds documented in the author's validation log were
conducted by separate AI agent instances as additional checks alongside the author's own verification. These checks
should not be interpreted as independent human expert review. The author designed the study, made or approved every
decision recorded in the frozen plan and validation log, approved the freezing of the study protocol, and takes full
responsibility for the paper and the archived materials.

\section*{Data and code availability}

The SAFE microdata are available from the ECB under its user agreement and are not redistributed. No record-level
SAFE data were linked with any other firm data. The code (extract preparation, card construction, elicitation, runner,
analysis and the checks of the frozen plan), the frozen plan with its development amendments, the frozen settings and
the aggregate reports are in the repository \url{https://github.com/speplinski/safe-profile-validation}, archived on
Zenodo (\url{https://doi.org/10.5281/zenodo.24396236}, all versions). The validation log, with every decision, review
round and correction, is kept in Polish in the author's project repository and is available from the author on request. Three levels of reproduction are distinct.
(i)~\emph{Reading the tables}: \texttt{scripts/safe\_paper\_numbers.py} prints every number of the tables from the
archived aggregate reports and checks that the confirmatory report was made with the archived frozen settings, plan
and analysis script; it reads the stored estimates and does not recompute them. (ii)~\emph{Recomputing from the
historical model answers}: the item-level answers are derived from SAFE records and are not published, so this needs
them from the author together with SAFE access. (iii)~\emph{A new replication}: \texttt{scripts/safe\_extract.py}
makes the Polish extract from the ECB archive \texttt{ecb.SAFE\_microdata.zip} (the version used has SHA-256
\texttt{64a118f1\ldots}, downloaded on 12 November 2026); with pandas 3.0.5 it reproduces the extract of this study
byte for byte (SHA-256 \texttt{15fef143\ldots}), from which \texttt{scripts/safe\_build.py} rebuilds the test items.
New model runs are then needed; served models behind API aliases may have changed and need not give the same answers,
and the analysis refuses runs on a serving stack other than the frozen one. A new replication should be analysed
under its own frozen settings, not reported as the original run. The confirmatory report did not record the versions
of the numerical libraries; the archived tests pass with Python 3.13, NumPy 2.5 and SciPy 1.18.

\section*{Acknowledgements and data source}

This paper uses data from the EC/ECB Survey on the access to finance of enterprises.

\begin{thebibliography}{99}

\bibitem[Aher et~al.(2023)]{aher2023}
G.~V. Aher, R.~I. Arriaga, and A.~T. Kalai.
\newblock Using large language models to simulate multiple humans and replicate human subject studies.
\newblock In \emph{Proceedings of the 40th International Conference on Machine Learning}, PMLR 202, pages 337--371, 2023.

\bibitem[Argyle et~al.(2023)]{argyle2023}
L.~P. Argyle, E.~C. Busby, N.~Fulda, J.~R. Gubler, C.~Rytting, and D.~Wingate.
\newblock Out of one, many: Using language models to simulate human samples.
\newblock \emph{Political Analysis}, 31(3):337--351, 2023.
\newblock doi:10.1017/pan.2023.2.

\bibitem[Bretz et~al.(2009)]{bretz2009}
F.~Bretz, W.~Maurer, W.~Brannath, and M.~Posch.
\newblock A graphical approach to sequentially rejective multiple test procedures.
\newblock \emph{Statistics in Medicine}, 28(4):586--604, 2009.
\newblock doi:10.1002/sim.3495.

\bibitem[Bretz et~al.(2011)]{bretz2011}
F.~Bretz, M.~Posch, E.~Glimm, F.~Klinglmueller, W.~Maurer, and K.~Rohmeyer.
\newblock Graphical approaches for multiple comparison procedures using weighted {B}onferroni, {S}imes, or parametric tests.
\newblock \emph{Biometrical Journal}, 53(6):894--913, 2011.
\newblock doi:10.1002/bimj.201000239.

\bibitem[Clopper and Pearson(1934)]{clopper1934}
C.~J. Clopper and E.~S. Pearson.
\newblock The use of confidence or fiducial limits illustrated in the case of the binomial.
\newblock \emph{Biometrika}, 26(4):404--413, 1934.
\newblock doi:10.1093/biomet/26.4.404.

\bibitem[Efron and Tibshirani(1993)]{efron1993}
B.~Efron and R.~J. Tibshirani.
\newblock \emph{An Introduction to the Bootstrap}.
\newblock Chapman \& Hall, 1993.

\bibitem[Epstein(1969)]{epstein1969}
E.~S. Epstein.
\newblock A scoring system for probability forecasts of ranked categories.
\newblock \emph{Journal of Applied Meteorology}, 8(6):985--987, 1969.
\newblock doi:10.1175/1520-0450(1969)008<0985:ASSFPF>2.0.CO;2.

\bibitem[Gneiting and Raftery(2007)]{gneiting2007}
T.~Gneiting and A.~E. Raftery.
\newblock Strictly proper scoring rules, prediction, and estimation.
\newblock \emph{Journal of the American Statistical Association}, 102(477):359--378, 2007.
\newblock doi:10.1198/016214506000001437.

\bibitem[Hall(1992)]{hall1992}
P.~Hall.
\newblock \emph{The Bootstrap and Edgeworth Expansion}.
\newblock Springer, New York, 1992.

\bibitem[Holm(1979)]{holm1979}
S.~Holm.
\newblock A simple sequentially rejective multiple test procedure.
\newblock \emph{Scandinavian Journal of Statistics}, 6(2):65--70, 1979.
\newblock \url{https://www.jstor.org/stable/4615733}.

\bibitem[Horton et~al.(2023)]{horton2023}
J.~J. Horton, A.~Filippas, and B.~S. Manning.
\newblock Large language models as simulated economic agents: What can we learn from {H}omo silicus?
\newblock NBER Working Paper 31122, 2023.
\newblock doi:10.3386/w31122.

\bibitem[Lakens(2017)]{lakens2017}
D.~Lakens.
\newblock Equivalence tests: A practical primer for t tests, correlations, and meta-analyses.
\newblock \emph{Social Psychological and Personality Science}, 8(4):355--362, 2017.
\newblock doi:10.1177/1948550617697177.

\bibitem[Santurkar et~al.(2023)]{santurkar2023}
S.~Santurkar, E.~Durmus, F.~Ladhak, C.~Lee, P.~Liang, and T.~Hashimoto.
\newblock Whose opinions do language models reflect?
\newblock In \emph{Proceedings of the 40th International Conference on Machine Learning}, PMLR 202, pages 29971--30004, 2023.

\end{thebibliography}

\end{document}
